% MTH5000 Masters Research Project
% Jaykumar Vinodbhai Chauhan, 34710280
%
% Every table in this document is generated by make_tables.py directly from the
% result files. Do not type a number in by hand: run make_tables.py instead.
%   python make_tables.py && latexmk -pdf main.tex
%
% Plain academic style: Times, black only, no colour, no decorative rules.

\documentclass[11pt,a4paper]{article}

\usepackage[T1]{fontenc}
\usepackage[utf8]{inputenc}
\usepackage{mathptmx}                 % Times for text and mathematics
\usepackage[margin=25mm]{geometry}
\usepackage{booktabs}
\usepackage{graphicx}
\usepackage{amsmath,amssymb}
\usepackage{caption}
\usepackage[hidelinks]{hyperref}
\usepackage{natbib}
\usepackage{setspace}

\captionsetup{font=small,labelfont=bf,justification=justified}
\setlength{\parskip}{0.4em}
\onehalfspacing

\title{\vspace{-1.5cm}\Large Machine Learning for Forecasting Hazardous Air
Quality Days in Indian Cities}

\author{%
Jaykumar Vinodbhai Chauhan \\ Student ID: 34710280 \\[0.4em]
Master of Mathematics \\
MTH5000 Masters Research Project \\[0.4em]
Supervisor: Dr Tianhai Tian \\
School of Mathematics, Monash University}

\date{\today}

\begin{document}
\maketitle

\begin{abstract}
\noindent
% TODO write last, once the discussion is settled. Roughly 200 words.
% It must state the threshold finding first, not the model comparison.
[Abstract to be written last. Lead with the decision threshold result, not with
which model won.]
\end{abstract}

% ============================================================================
\section{Introduction}

Air pollution in Indian cities regularly reaches concentrations classified as
hazardous to human health. The public response to such episodes depends on
knowing that a severe day is coming before it arrives: a forecast issued on the
morning of a hazardous day is of limited value, while one issued a day or more
ahead can change behaviour.

This project treats that as a forecasting problem and, crucially, evaluates it as
a warning system rather than as a curve-fitting exercise. Three research
questions follow.

My interest in this question is also personal. Having grown up in India, air
pollution was not an abstract statistic but part of ordinary life: seasonal
smog that turned daylight hazy, and days when going outside meant a sore
throat. Whether those days can be predicted reliably enough to be worth
warning people about is a question I have a genuine stake in.

\begin{enumerate}
\item How accurately can daily concentrations be forecast one to three days
ahead by machine learning, relative to classical and naive benchmarks?
\item Is it better to forecast the concentration and then apply a threshold, or
to classify the exceedance event directly?
\item Considered as a warning system, what hit rate and false alarm rate does
each approach achieve, and how does the decision threshold trade them off?
\end{enumerate}

The answer to the third question turns out to dominate the other two, and the
report is organised accordingly.

% ============================================================================
\section{Literature}
% TODO Jay. This is the one section I cannot write for you, because every
% citation must be verified against CrossRef before it goes in and I will not
% invent one. Four things to cover, with a suggested search for each:
%
%   1. Air quality forecasting in Delhi and northern India. Search: "PM2.5
%      forecasting Delhi machine learning", "stubble burning air quality Delhi".
%   2. Machine learning versus classical methods on time series. The M4 and M5
%      forecasting competition papers are the standard references for the claim
%      that simple methods are hard to beat on limited data.
%   3. Warning system verification. Search: "probability of detection false
%      alarm rate forecast verification", "critical success index meteorology".
%      Jolliffe and Stephenson is the standard text.
%   4. Negative and invalid readings in reference monitors. There is a 2023
%      paper, "On thresholds for controlling negative particle (PM2.5) readings
%      in air quality reporting", directly on point. Verify it.

[To be written. See the notes in the source file for the four areas to cover and
the searches to run. Verify every citation against CrossRef before use.]

% ============================================================================
\section{Data}
\label{sec:data}

\subsection{Source and station selection}

Daily mean PM$_{2.5}$ concentrations were obtained from OpenAQ. Five candidate
monitoring stations within 25\,km of central Delhi were compared on data
coverage rather than on nominal length of record.

The comparison produced an unexpected constraint. Every station operated by the
Central Pollution Control Board returns measurements only from 19 February 2025,
approximately eighteen months, although the OpenAQ location metadata advertises a
first observation in February 2016. Four independent stations sharing an
identical cutoff indicates a systematic limit on what the measurements endpoint
serves rather than a coincidence. Only location 8118 returns the full record and
it was therefore selected.

This carries a limitation that must be stated plainly. Location 8118 is supplied
through AirNow and is a United States diplomatic post monitor rather than part of
the Indian official network. The project therefore forecasts at a single
well-characterised reference site in Delhi, which is common in this literature,
and the results should not be read as representative of the city as a whole.

\begin{table}[htbp]
\centering
\caption{The analysis dataset after cleaning.}
\label{tab:dataset}
\input{tables/dataset}
\end{table}

\subsection{Two defects in the source data}
\label{sec:defects}

Two problems were identified before any modelling, one of which was silently
corrupting the series.

\paragraph{Physically impossible values.} The raw series contained five negative
concentrations and a maximum of 1990\,\textmu g\,m$^{-3}$. The value 1990.0
appeared three times identically on unrelated dates, which is the signature of a
clipped instrument ceiling rather than a measurement. One occurrence fell on 2
September, during the monsoon, between neighbouring days of 42 and 178.

\paragraph{An incorrect coverage denominator.} OpenAQ reports the completeness of
each daily aggregate as the observation count divided by an expected count of 24,
assuming hourly reporting. The sensor at location 8118 reported half-hourly from
2016 to 2024 and hourly from 2025 onward. The visible symptom is a reported
completeness exceeding 100 per cent, which is impossible. The damaging symptom is
invisible: a day holding 24 of a genuine 48 readings is scored as fully complete,
so a mean computed from half a day of a polluted afternoon enters the series as
though it were a daily mean. Days at 50 to 75 per cent true coverage produced
implausible values approximately seven times as often as days above 90 per cent,
which identifies incomplete days as the mechanism rather than a coincidence.

The correction takes the reporting cadence from the data, using the ninetieth
percentile of observation counts within each year, which follows the 2025 change
in reporting interval automatically. Recomputed on that basis, 1{,}085 of 3{,}199
observed days fall below 75 per cent coverage rather than the 268 that the
published completeness figure implies.

\begin{table}[htbp]
\centering
\caption{Cleaning rules applied, and the basis for each.}
\label{tab:cleaning}
\input{tables/cleaning}
\end{table}

The coverage threshold is set at 50 per cent rather than a stricter value because
the longest gap is the binding constraint: at 60 per cent the worst gap rises to
69 days and at 90 per cent to 185 days, either of which would distort estimation
of the annual cycle. At 50 per cent the worst gap is 34 days. The exceedance rate
moves only from 30.0 to 27.0 per cent across every candidate threshold, which
confirms that the filter is not selectively removing the events of interest.

\begin{figure}[htbp]
\centering
\includegraphics[width=\textwidth]{figures/fig2_coverage}
\caption{The coverage defect. Panel (a) shows the number of readings contributing
to each daily mean, with the cadence estimated from the data each year as the
solid line and the fixed denominator of 24 as the dashed line. Panel (b) shows
the rate of implausible values by true coverage.}
\label{fig:coverage}
\end{figure}

\begin{figure}[htbp]
\centering
\includegraphics[width=\textwidth]{figures/fig1_series}
\caption{The cleaned daily series, with days exceeding the threshold shaded. The
episodic structure is the feature the project sets out to forecast: exceedances
arrive as multi-day winter episodes rather than as isolated days.}
\label{fig:series}
\end{figure}

% ============================================================================
\section{Methodology}
\label{sec:method}

\subsection{An assumption that the data did not support}
\label{sec:weekly}

The approved proposal states that daily pollution series carry a weekly cycle
driven by traffic and industrial activity, and treats the coexistence of a weekly
and an annual period as the central modelling difficulty, since a seasonal ARIMA
accommodates only one seasonal period.

Tested on the cleaned series, the weekly cycle is not present at this station.
Three independent tests agree. The periodogram gives dominant periods of 357.9,
178.9, 188.4, 397.7, 325.4 and 60.7 days, which is the annual cycle and its
harmonics, with nothing near seven. On the deseasonalised logarithm of
concentration the autocorrelation at lag 7 is 0.173, while lag 6 is 0.179 and lag
8 is 0.153, so lag 7 lies on the decay curve rather than above it; the same holds
at lags 14 and 21. The day of week effect on the same residuals spans 6.5 per
cent between the highest and lowest day, with an analysis of variance giving
$F = 1.13$, $p = 0.34$ and a Kruskal-Wallis test giving $p = 0.77$.

\begin{figure}[htbp]
\centering
\includegraphics[width=\textwidth]{figures/fig3_weekly}
\caption{No weekly cycle. Lags 7, 14 and 21 are circled and lie on the decay
curve rather than above it. Every day of week interval contains zero.}
\label{fig:weekly}
\end{figure}

There is a physical reading rather than a null result. A weekly rhythm in urban
particulate concentration is a traffic signature, and location 8118 sits in the
diplomatic enclave rather than at a roadside site such as ITO or Anand Vihar. Its
absence at a site that is not traffic dominated is what one would expect.

The consequence for the specification is direct: the classical benchmark is ARIMA
with annual Fourier terms as exogenous regressors and no weekly seasonal order.
The specification is simpler, better identified, and empirically justified rather
than assumed.

\begin{figure}[htbp]
\centering
\includegraphics[width=0.66\textwidth]{figures/fig4_annual}
\caption{The annual cycle, by contrast, is unmistakable: monthly median with the
interquartile range shaded. Its strength is what makes the absence in
Figure~\ref{fig:weekly} credible rather than an artefact of a weak test.}
\label{fig:annual}
\end{figure}

\subsection{Feature construction}

A supervised table was built with one row per forecast origin. A row indexed by
date $t$ carries only functions of observations up to and including day $t$, and
its targets are the concentrations on days $t+1$, $t+2$ and $t+3$. Features
comprise lags at 0, 1, 2, 3 and 7 days; first differences; rolling mean, standard
deviation and maximum over 3, 7, 14 and 30 days; an anomaly and standardised
anomaly against the 30 day norm; exceedance frequency over the last 7 and 30
days; days since the last exceedance; calendar terms; four pairs of annual
Fourier terms; and three missingness indicators.

Missing values are never imputed. Gradient boosting accepts them natively and the
state space models skip the measurement update at a gap, so the absence is
carried as information rather than replaced by a fabricated value. The single
exception is the recurrent network, discussed in Section~\ref{sec:gru}.

\subsection{Models}

Ten specifications were fitted. Three are benchmarks: persistence, the seasonal
naive forecast at lag 7, and a smoothed day of year climatology whose smoothing
window was selected on a validation block. The classical model is ARIMA with
annual Fourier terms as exogenous regressors, fitted on the logarithm of
concentration, with the order selected from an explicit grid by the Akaike
information criterion on the training block. Four are machine learning models
fitted to the feature table: penalised linear regression, and random forest and
histogram gradient boosting each fitted both to the level and to the change from
the origin. The last is a gated recurrent network.

Working in logarithms is not cosmetic. Annual Fourier terms explain 69 per cent
of the variance of the logarithm against considerably less on the raw scale, and
fitting in logs prevents the model predicting negative concentrations.
Back-transformation is not innocent either: if $\log y$ has mean $m$ and variance
$s^2$ then $\exp(m)$ is the median of $y$ and $\exp(m + s^2/2)$ is its mean. Both
were produced as separately named models, since a report that back-transforms
without stating which quantity it reports has silently changed what is being
forecast.

\subsection{Rolling origin evaluation}
\label{sec:rolling}

Every model was evaluated over 2{,}303 origins from 9 November 2019 to 24 August
2026, spanning 6.8 years, with a three year burn-in, parameters re-estimated
every 90 days and model state advanced daily in between.

This is not a refinement of a single train-test split; it replaces one. On a
single test block the ARIMA scored 0.940 in relative terms at one day ahead
against 0.903 on the single block, because that block happened to be calmer than
the training period and was flattering every model. A single split is one draw
from a process, and a conclusion drawn from it is a conclusion about that period.

The governing invariant is that at origin $t$ a model may use observations up to
and including day $t$ and nothing else: not the value at $t+1$, not a parameter
estimated using data after $t$, and not a scaling constant or imputed value
computed from the full series. Since nothing fails visibly when that is broken,
and the numbers merely improve, the invariant is enforced by construction and
tested directly. Every model runs through one shared loop, and that loop is
verified by corrupting the series from a chosen date, rerunning, and requiring
that no forecast issued before that date changes.

One trap is specific to supervised learning under a rolling origin. At origin $t$
the label attached to feature row $t'$ is the concentration at $t' + h$, so the
training set ends at $t - h$ rather than at $t$. A separate model is fitted per
horizon because that arithmetic requires it. The error would have amounted to
three rows per refit and would have been entirely invisible in the output.

\subsection{The recurrent network and what it cannot do}
\label{sec:gru}

The recurrent model is a single-layer gated recurrent unit over a 30 day window,
each day carrying the logarithm of concentration, a missingness flag and annual
Fourier terms, with 24 hidden units and early stopping on a chronological holdout
inside the training window. It predicts the change in log concentration from the
origin, with the origin value added back afterwards.

It also forces a compromise that no other model in the suite requires. A
recurrent network must be handed a value at every timestep: it cannot skip one
the way the Kalman filter skips a measurement update, and it cannot accept a
missing value the way gradient boosting can. Gaps are therefore carried forward
with a separate channel marking which days were filled. The architecture the
project was asked to include is the only one in the suite that cannot represent
the statement that a value is unknown.

\subsection{Evaluation measures}

Accuracy is reported separately at each horizon and never averaged, since skill
decays with horizon and concealing that would defeat the purpose. The headline
measure is relative mean absolute error against the persistence forecast computed
on identical rows, so that 1.000 is persistence and values below one beat it.
Mean absolute scaled error is reported alongside, but it divides by the training
period volatility and therefore also reflects any difference in volatility
between periods, which is material here.

As a warning system, a day is an event if the observed concentration exceeds the
health threshold, and an alarm is raised when the forecast exceeds a decision
threshold. These are separate quantities and are kept separate throughout.
Reported are the hit rate, the false alarm rate defined as false alarms divided
by quiet days, precision, the critical success index, and alarms raised per year.
The last is included because a false alarm rate of 0.10 sounds negligible and
means roughly one wasted alarm every ten quiet days, which is what an operator
actually experiences.

% ============================================================================
\section{Results}
\label{sec:results}

\subsection{Forecast accuracy}

\begin{table}[htbp]
\centering
\caption{Relative mean absolute error against persistence, under rolling origin
evaluation over 2{,}303 origins. Lower is better and 1.000 is persistence. The
recurrent network is given as the range over five initialisation seeds; it is the
only model whose result depends on one.}
\label{tab:accuracy}
\input{tables/accuracy}
\end{table}

Table~\ref{tab:accuracy} answers the first research question, and the answer is
not uniform across horizons.

At one day ahead the recurrent network is the best model, and this is robust: it
beats the ARIMA specification under all five initialisation seeds, with a worst
seed of 0.935 against 0.940. At two days the ARIMA is better under all five
seeds, though by 0.003 against the best of them. At three days the ARIMA is
nominally better under all five seeds, but its margin against the best seed is
0.0003, which is a tie in any language that would survive examination.

The mechanism is coherent. At one day ahead the forecast is dominated by
short-range structure, where a nonlinear model has something to exploit. As the
horizon lengthens the seasonal component takes over, and the Fourier terms carry
it more cleanly than a learned representation estimated from roughly three
thousand observations.

The network is reported as a range rather than a point because it is the only
model whose result depends on a random initialisation, and the dependence is not
small. Across seeds the spread is 0.019 at one day, 0.013 at two and 0.047 at
three. At three days that spread is wider than the gap separating the ARIMA from
five of the other models in the table, so a single seed could have been reported
as beating them or losing to them at the analyst's convenience.

Two subsidiary results support the specification choices. Removing the Fourier
terms from the ARIMA degrades it by 4.5 points of relative error at one day and
by 10 points at three, so the seasonal terms earn their place and increasingly so
with horizon. Modelling the change from the origin rather than the level improves
both tree models, because a tree cannot extrapolate beyond the targets it was
trained on and this series drifts downward across the decade.

\subsection{Two arms of the same question}

% TODO Jay: expand this once you have decided how much of the classifier detail
% to report. The numbers are in threshold_sweep.csv.

Neither arm dominates. Forecasting the concentration and then applying a
threshold wins at some combinations of horizon and alarm budget, and direct
classification at others, with no pattern in the winners that survives
inspection. Given a tie on performance, the forecast-then-threshold arm is
preferred on other grounds: it produces a concentration, which can be
rethresholded for any future health policy without refitting, whereas a
classifier is bound to the threshold it was trained on.

\subsection{The decision threshold}
\label{sec:threshold}

\begin{table}[htbp]
\centering
\caption{Best hit rate available within a budget of alarms per year, with the
event fixed at the health threshold and only the decision threshold varying.}
\label{tab:budget}
\input{tables/budget}
\end{table}

Table~\ref{tab:budget} contains the principal result of this project. Moving the
alarm budget from 40 to 100 per year changes the achievable hit rate at one day
ahead from 0.377 to 0.873. Holding the budget fixed and changing the model
instead moves it by around five points.

\begin{figure}[htbp]
\centering
\includegraphics[width=\textwidth]{figures/fig5_warning}
\caption{The trade-off, by horizon. The top row shows hit rate against false
alarm rate; the bottom row replaces the false alarm rate with alarms raised per
year, the unit in which an operator budgets. The curves very nearly coincide,
which is the finding. They have deliberately not been separated for legibility.}
\label{fig:warning}
\end{figure}

Two individual comparisons make the point sharper than the aggregate does. At a
budget of 100 alarms per year and one day ahead, persistence achieves a hit rate
of 0.873, which is the best figure any model in the study achieves at that budget.
At three days ahead and a budget of 60, a climatological forecast, which uses no
recent information whatever, matches the best model to within four thousandths.
Figure~\ref{fig:warning} shows the same result as curves; they very nearly
coincide, and they have deliberately not been separated for legibility.

The reading is not that modelling is pointless. It is that the quantity a
decision maker controls, namely how many alarms per year can be justified,
dominates the quantity a modeller controls. A report that presents a model
comparison without stating the operating point it was evaluated at has answered a
question nobody asked.

% ============================================================================
\section{Discussion}
\label{sec:discussion}

\subsection{The operational question dominates the modelling question}

Table~\ref{tab:budget} and Figure~\ref{fig:warning} together make a single
point, and this section states what follows from it rather than repeating the
numbers. A decision maker who fixes how many alarms per year an operator or
the public will tolerate has more influence over how many hazardous days are
caught than a modeller choosing among the ten specifications fitted in this
project. Moving the budget from 40 to 100 alarms a year moves the achievable
hit rate at one day ahead from 0.377 to 0.873, a range of fifty points.
Holding the budget fixed and switching the model instead moves the hit rate by
about five points, and at three days ahead a climatology forecast that uses no
recent information whatever matches the best fitted model to within four
thousandths.

This is not an argument against modelling. ARIMA and the recurrent network
both earn real, if modest, gains over persistence in Table~\ref{tab:accuracy},
and a better point forecast still narrows the range of thresholds worth
considering. It is an argument about where the leverage in this problem
actually sits. A hazardous air quality warning system is not a machine
learning benchmark with a threshold attached to it afterwards; the threshold
is the policy instrument, and the forecasting model only supplies its input.
A comparison of models that does not first fix the alarm budget it was
evaluated at has answered a question nobody who has to operate the system
actually asked. The practical recommendation that follows is to treat the
alarm budget as the primary subject of discussion with the health authority
commissioning such a system, and the choice of forecasting model as a
secondary refinement of a curve whose shape is already mostly fixed by what
every model in this study shares: the persistence of concentration from one
day to the next, and the annual cycle.

\subsection{How large is a difference between models?}

Two results in this project bound that question from below. Running the identical
pipeline on two different machines moved the tree models by up to 0.004 in
relative error, because histogram gradient boosting parallelises with OpenMP and
floating point summation order depends on the thread count. Changing only the
random initialisation of the recurrent network moved it by 0.013 to 0.047.

Several of the gaps between adjacent models in Table~\ref{tab:accuracy} are of
that order. Any claim that one of those models is better than another, made
without stating the run to run variation, is not supported by the evidence.

\subsection{Limitations}

The analysis rests on a single monitor, OpenAQ location 8118, which
Section~\ref{sec:data} shows is supplied through AirNow and is almost
certainly a United States diplomatic post instrument rather than a station in
the Central Pollution Control Board network. It was chosen because it is the
only nearby monitor with a complete decade of record, against the eighteen
months every CPCB station returns; the choice buys history at the cost of
representativeness. The absence of a weekly cycle documented in
Section~\ref{sec:weekly} is consistent with a diplomatic enclave rather than a
roadside site, which supports this reading of the limitation rather than
leaving it an unrelated coincidence, but it does not change the fact that
every number in this report describes one point in the city rather than Delhi
as a whole.

The project forecasts one city. Delhi was chosen because its air quality
problem is severe and well documented, not because it is representative of
Indian cities generally. Pollution sources, meteorology, and the strength of
the seasonal cycle that drives the Fourier terms in Section~\ref{sec:weekly}
all vary by region, and a coastal city or one without a dominant biomass
burning season could show a different balance between the classical and
machine learning models, or a genuine weekly cycle where this station shows
none. Nothing in this report should be read as a claim about any station or
city other than the one it was fitted to.

The upper cleaning threshold of 1000\,\textmu g\,m$^{-3}$ is a judgement rather
than a fact, and its effect on the reported conclusions was checked directly
rather than assumed. Restoring the seven days it removes, including the three
identical readings at 1990 diagnosed in Section~\ref{sec:defects} as a clipped
instrument ceiling, gives 3{,}100 usable days and 927 exceedance days rather
than 3{,}093 and 920, a small shift. Refitting the full rolling-origin harness
on this uncapped series leaves the ARIMA specification essentially unchanged,
0.935 against 0.940 at one day ahead and 0.829 against 0.832 at three days,
because working in logarithms and differencing the series absorbs the
contamination. The machine learning models are not so protected. Ridge moves
from a near tie with persistence to a clear loss at one day ahead, 0.997 to
1.118, and gradient boosting on the level target moves from a narrow win over
persistence to a loss at three days, 0.883 to 1.006. The central finding of
Research Question 1, that the classical benchmark beats every machine learning
model at every horizon, therefore does not merely survive removing the
threshold; it strengthens, because the models it beat most narrowly are
precisely the ones most exposed to the uncleaned tail. The threshold matters,
but it matters more to the machine learning suite than to the model this
project ultimately recommends.

Hyperparameters for the tree-based models were chosen at every refit from a
small explicit grid, scored on a held-out slice of the training window taken
in time order, rather than from an exhaustive or Bayesian search. This
followed a real failure during development, in which fixed hyperparameters let
gradient boosting overfit badly enough to score worse than persistence; the
grid search fixed that specific problem, but a small grid can still miss a
configuration that a wider search would find. The machine learning results in
Table~\ref{tab:accuracy} should therefore be read as a lower bound on what
those model classes could achieve with more computation, not as their ceiling.

No meteorological covariate enters any model. Wind speed in particular is a
well established driver of particulate dispersion and is a plausible source
of the short-range structure that lets the recurrent network beat ARIMA at one
day ahead. A version of this project with wind, temperature, and boundary
layer height added to the feature table is the most direct extension; the
\texttt{meteostat} package, which wraps NOAA and GHCN station data, would
supply it without a new data acquisition problem. Its absence here is a scope
decision made to fit the project inside nine weeks, not a claim that
meteorology is unimportant.

% ============================================================================
\section{Conclusion}

This project set out to forecast hazardous air quality days in an Indian city
and to evaluate the forecasts as a warning system rather than as a
curve-fitting exercise, and the answer to the third research question turned
out to dominate the other two. Across ten models spanning simple benchmarks, a
classical time series model, four machine learning regressors, and a recurrent
network, the spread in achievable hit rate at a fixed alarm budget is about
five percentage points; the spread across alarm budgets, holding the model
fixed, is about fifty. The choice a decision maker makes about how many alarms
per year can be justified therefore has far more influence on how many
hazardous days are caught than the choice a modeller makes among the
specifications fitted here.

Machine learning earned its place only partially. ARIMA with annual Fourier
terms beats every machine learning regressor at every horizon, a result that
strengthens rather than weakens under the sensitivity analysis in
Section~\ref{sec:discussion}. The one machine learning method that does win,
the recurrent network, does so only at one day ahead, and only when reported
honestly as a range across five random initialisations rather than as a single
favourable seed. The supervisor's instruction to try several methods rather
than defaulting to ARIMA was the right one: it is what makes the negative
result at two and three days credible, and what surfaces the one day ahead
result that a project restricted to classical methods would never have found.

Two extensions follow directly from the limitations above. Adding
meteorological covariates, wind speed in particular, is the most direct route
to improving the short horizon forecasts that the recurrent network already
shows have room to improve. Repeating the analysis at a second, independent
station, ideally one inside the Central Pollution Control Board network once
its record lengthens, would test whether the findings here are a property of
Delhi air quality or an artefact of the one diplomatic enclave monitor this
project was restricted to. Until then, the central recommendation stands on
its own: a hazardous air quality warning system should be designed around the
alarm budget its operator can justify, with the forecasting model treated as a
secondary refinement rather than the primary design choice.

% ============================================================================
\section*{Reproducibility}

The analysis runs end to end from an empty directory in eight command line
scripts, each carrying an automated correctness test that runs before it produces
any output. The rolling origin harness verifies its own leakage invariant before
every run. Every table in this report is generated directly from the result files
by \texttt{make\_tables.py}; no number has been transcribed by hand. Raw and
cleaned data are held under version control so that the analysis does not depend
on the data provider serving identical values at a later date, which the finding
in Section~\ref{sec:data} suggests is not a safe assumption.

% ============================================================================
% \bibliographystyle{plainnat}
% \bibliography{refs}

\end{document}
